%% ma: L12-B08-0098
%% lop: 12
%% chuong: 2
%% bai: 8
%% dang: 12.8.29
%% loai: TN4
%% muc: NB
%% dap_an: "B"
%% nguon: "CĐ2.3.pdf giáo viên gửi 27/10/2026 (M07, Lê Minh Khôi – Nguyễn Anh Khôi, Chương II Chủ đề 3 Biểu thức tọa độ của các phép toán vectơ), Câu 7"
%% kien_thuc: "Bài 8 (biểu thức tọa độ của phép cộng, trừ vectơ, tích vectơ với một số)"
%% trang_thai: chinh-thuc
\begin{ex}
Trong không gian $Oxyz$, cho ba vectơ $\vec u=(2;-5;3)$, $\vec v=(0;2;-1)$ và $\vec w=(1;7;2)$. Tọa độ của vectơ $\vec a=\vec u-4\cdot\vec v-2\cdot\vec w$ là
\choice{$\vec a=(4;1;11)$}{\True $\vec a=(0;-27;3)$}{$\vec a=(0;1;3)$}{$\vec a=(4;-27;11)$}
\loigiai{$\vec a=(2-0-2;-5-8-14;3+4-4)=(0;-27;3)$.}
\end{ex}
